% ==================================================================
%  Preliminaries.tex  --  Section 2
%  The model setting moved to the Technical Overview (Section 1);
%  its label sec:model is defined there.
%  Cross-references used here:
%    \label{sec:method}    FHE-DBSCAN (proposed method) section
%    \label{sec:intro}     Introduction (Section 1)
% ==================================================================
\section{Preliminaries}
\label{sec:prelim}

\subsection{Notation}
\label{sec:notation}

Table~\ref{tab:notation} lists the notation used throughout the paper.

\begin{table}[t]
\centering
\caption{Notation.}
\label{tab:notation}
\renewcommand{\arraystretch}{1.15}
\begin{tabular}{@{}l p{0.75\linewidth}@{}}
\toprule
Symbol & Description \tabularnewline
\midrule
$X$
  & \raggedright dataset $\{\mathbf{x}_0, \dots, \mathbf{x}_{n-1}\} \subset
    \mathbb{R}^d$ \tabularnewline
$n$, $d$ & \raggedright number of points, dimension \tabularnewline
$\varepsilon$, $\mathit{MinPts}$
  & \raggedright DBSCAN radius and density threshold \tabularnewline
$\varepsilon_n$
  & \raggedright radius after the normalization of
    Section~\ref{sec:client} \tabularnewline
$N$ & \raggedright ring degree of CKKS; a ciphertext has $N/2$ slots
  \tabularnewline
$\mathsf{ct}.\mathbf{x}$
  & \raggedright ciphertext that encrypts the vector $\mathbf{x}$
    \tabularnewline
$\mathsf{ct}_j$
  & \raggedright ciphertext of the $j$-th coordinates of all $n$ points,
    $j = 1, \dots, d$ \tabularnewline
$\boldsymbol{\ell}$, $\mathsf{adj}_k$
  & \raggedright in Section~\ref{sec:method}, a bold or sans-serif vector
    symbol denotes the ciphertext of that vector, whose slot $i$ carries
    the value of the $i$-th point \tabularnewline
$\mathbf{v}[i]$
  & \raggedright slot $i$ of $\mathbf{v}$; the client puts the $i$-th
    point of the order of Section~\ref{sec:client} in slot $i$,
    $i = 0, \dots, n-1$ \tabularnewline
$\mathsf{Rot}(\mathbf{v}, k)$
  & \raggedright rotation of all $N/2$ slots: slot $i$ holds
    $\mathbf{v}[(i+k) \bmod N/2]$ \tabularnewline
$\mathsf{Rot}_k(\mathbf{v})$
  & \raggedright cyclic shift of the first $n$ slots: slot $i$ holds
    $\mathbf{v}[(i+k) \bmod n]$; realized with two rotations and two
    plaintext masks \tabularnewline
$\mathsf{Sgn}$
  & \raggedright composite minimax approximation of $\operatorname{sgn}$,
    applied to every slot; its precision $\alpha$ is set per stage
    (Section~\ref{sec:correct}) \tabularnewline
mask vector
  & \raggedright vector whose slots are $0$ or $1$; multiplying by it
    with $\odot$ keeps the slots where it is $1$ and sets the others to
    $0$ \tabularnewline
$k_{\max}$
  & \raggedright largest slot offset the server compares
    (Section~\ref{sec:client}) \tabularnewline
$R$
  & \raggedright number of propagation rounds (Section~\ref{sec:client})
    \tabularnewline
$\mathsf{adj}_k$
  & \raggedright adjacency mask for offset $k$
    (Section~\ref{sec:neighbor}) \tabularnewline
$\mathsf{Core\_mask}$
  & \raggedright core mask (Section~\ref{sec:core}) \tabularnewline
$\boldsymbol{\ell}$
  & \raggedright label vector (Section~\ref{sec:lp}) \tabularnewline
\bottomrule
\end{tabular}
\end{table}

\subsection{The CKKS Scheme}
\label{sec:ckks}

CKKS~\cite{cheon2017homomorphic} is a homomorphic encryption scheme for
approximate arithmetic over the real and complex numbers. It is built on
the ring
$R = \mathbb{Z}[X]/(X^{N}+1)$, where the degree $N$ is a power of two, and
on its residue rings $R_{q} = R/qR$. A message is a vector of up to $N/2$
complex numbers: the $N$ complex roots of $X^{N}+1$ come in $N/2$
conjugate pairs, so a polynomial with real coefficients carries only
$N/2$ independent complex values. We use real vectors throughout.
Encoding carries such a vector into $R$ through the inverse of the
canonical embedding, scaled by a factor $\Delta$ and rounded to integer
coefficients. The encoded message is therefore a fixed-point
representation of the vector: $\Delta$ fixes how many fractional bits are
kept, and the rounding discards everything below them. A ciphertext at
level $\ell$ is a pair of polynomials in $R_{q_{\ell}}$, where $q_{\ell}$
is the $\ell$-th modulus of a chain
$q_{0} \mid q_{1} \mid \dots \mid q_{L}$, and a fresh ciphertext starts at
the top level~$L$. We use the residue number system (RNS) variant of
CKKS~\cite{cheon2018full}, which DesiloFHE implements: each $q_{\ell}$ is
a product of word-sized primes, rescaling to the level below drops the
last of them, and every ring element is stored through the Chinese
remainder theorem as its residues modulo these primes. Security rests on
the ring learning with errors assumption. Each entry of the message occupies one \emph{slot};
additions and multiplications act on all $N/2$ slots at once, in a
single-instruction-multiple-data (SIMD) manner, which is what lets one
ciphertext carry a whole dataset. Decryption does not return the message
exactly: the rounding of the encoding and the noise of the encryption sit
below the fixed-point bits and grow as operations are chained, so every
value is only correct to a precision that the design has to keep track
of. Throughout this paper $N$ denotes the ring degree and $n$ the number
of data points.

We use the algorithms below, each described by what its output
encrypts.

\begin{itemize}
\item $\mathsf{KeyGen}$ produces the secret key $\mathsf{sk}$, the public
  key $\mathsf{pk}$, and the evaluation keys $\mathsf{evk}$. The
  evaluation keys are a relinearization key, a rotation key for the
  shifts that occur in the computation, a conjugation key, and a
  bootstrapping key. They allow the operations below to be carried out
  without $\mathsf{sk}$.
\item $\mathsf{Enc}_{\mathsf{pk}}(\mathbf{x})$ encodes $\mathbf{x}$ with
  the scaling factor $\Delta$ and returns a ciphertext
  $\mathsf{ct}.\mathbf{x}$ of $\mathbf{x}$;
  $\mathsf{Dec}_{\mathsf{sk}}(\mathsf{ct}.\mathbf{x})$ returns an
  approximation of $\mathbf{x}$.
\item $\mathsf{Add}(\mathsf{ct}.\mathbf{x}, \mathsf{ct}.\mathbf{y})$
  returns an encryption of $\mathbf{x} + \mathbf{y}$, and
  $\mathsf{Sub}(\mathsf{ct}.\mathbf{x}, \mathsf{ct}.\mathbf{y})$ returns
  an encryption of $\mathbf{x} - \mathbf{y}$. One of the two operands may
  be a plaintext vector. Neither operation consumes a level.
\item $\mathsf{Mul}(\mathsf{ct}.\mathbf{x}, \mathsf{ct}.\mathbf{y})$
  returns an encryption of $\mathbf{x} \odot \mathbf{y}$, the slot-wise
  (Hadamard) product. The product of two ciphertexts consists of three
  polynomials instead of two. Relinearization with the relinearization
  key brings it back to two, and rescaling then divides the scaling factor
  back to $\Delta$ and moves the ciphertext one level down the modulus
  chain. Multiplying by a plaintext vector encoded with the scaling factor
  $\Delta$, which we use for masks and scalars, needs no relinearization,
  but the result is also rescaled and consumes one level.
\item $\mathsf{Rot}(\mathsf{ct}.\mathbf{x}, k)$ returns an encryption of
  the vector whose $i$-th slot is $x_{i+k}$, the indices being taken
  modulo the number of slots. It is a Galois automorphism followed by key
  switching under the rotation key, and consumes no level.
  $\mathsf{Conj}$ conjugates the slots and is used inside bootstrapping.
\end{itemize}

Once a ciphertext has spent its levels, bootstrapping with the
bootstrapping key returns an encryption of the same vector at a higher
level. It is by far the most expensive operation of the scheme, and the
number of bootstrappings, rather than the number of additions or
multiplications, sets the running time of our pipeline. Because it
changes no slot value, we leave it out of the algorithms of
Section~\ref{sec:method} and apply it wherever a ciphertext runs out of
levels.

The algorithms of Section~\ref{sec:method} are expressions on
ciphertexts. There we name a ciphertext after the vector it encrypts,
such as $\mathsf{adj}_k$ or $\boldsymbol{\ell}$, and write the
operations above as $\mathbf{x} + \mathbf{y}$, $\mathbf{x} - \mathbf{y}$,
$\mathbf{x} \odot \mathbf{y}$, $c\,\mathbf{x}$ for a public scalar $c$,
and $\mathbf{x}^{2}$ for $\mathbf{x} \odot \mathbf{x}$. Rotations appear
there mostly as the cyclic shift $\mathsf{Rot}_k$ of
Table~\ref{tab:notation}, which is built from two calls to
$\mathsf{Rot}$. Label Propagation also uses $\mathsf{Rot}$ itself, to
move blocks of $n$ slots within the ciphertext.

Comparison, division, and any other non-polynomial function are not part
of this set. Everything a DBSCAN evaluation needs beyond additions,
multiplications and rotations must therefore be built from polynomials;
Section~\ref{sec:comparison} does this for the comparison.

\subsection{DBSCAN}
\label{sec:dbscan}

DBSCAN~\cite{ester1996density} defines clusters by local density rather
than by distance to a representative point. Given a radius $\varepsilon$
and a threshold $\mathit{MinPts}$, the $\varepsilon$-neighborhood of a
point $p$ in a dataset $X$ is
$\mathcal{N}_{\varepsilon}(p) = \{\, q \in X : \mathrm{dist}(p,q) \le \varepsilon
\,\}$, which contains $p$ itself, and $p$ is a \emph{core point} if
$|\mathcal{N}_{\varepsilon}(p)| \ge \mathit{MinPts}$. In one sentence, the two
parameters fix a single density threshold: a region is dense enough to
carry a cluster when each of its points has at least $\mathit{MinPts}$
points within distance $\varepsilon$.

A point $q$ is \emph{directly density-reachable} from a point $p$ if $p$
is a core point and $q \in \mathcal{N}_{\varepsilon}(p)$. A point $q$ is
\emph{density-reachable} from $p$ if there is a chain of points
$p = p_1, p_2, \dots, p_m = q$ in which each $p_{j+1}$ is directly
density-reachable from $p_j$; every point of the chain except possibly
$q$ is then a core point. Two points are \emph{density-connected} if both
are density-reachable from a common point, and a cluster is a maximal set
of density-connected points. Intuitively, a cluster grows from a core
point by repeatedly adding the neighbors of its core points. A non-core
point that lies in the neighborhood of a core point is a \emph{border
point} and joins that cluster; a point that is neither core nor border
belongs to no cluster and is reported as noise.

Two consequences matter for this paper. First, because membership follows
chains of local links, a cluster can take an arbitrary shape---elongated,
curved, or wrapped around another cluster---and the number of clusters
need not be known in advance. This is what separates DBSCAN from
K-means (Section~\ref{sec:intro}), whose clusters are the cells of a set
of centroids and are therefore convex.

Second, border points are ambiguous. A border point can lie in the
neighborhoods of core points of two different clusters, and the
definition does not say which one it belongs to. Implementations resolve
this in several ways. The original algorithm follows the order of the
traversal, so that the cluster whose expansion reaches the point first
takes it. Parallel algorithms built on a disjoint-set (union--find)
structure~\cite{patwary2012new} merge core points into clusters and
attach each border point to only one of them. A fixed rule, such as
taking the cluster with the largest label, removes the choice altogether.
Each of these yields a valid DBSCAN output. We adopt the maximum rule, because it is the one that can be evaluated on
ciphertexts: it is a maximum over encrypted labels, whereas a first-come
rule presupposes a traversal whose order depends on the data, which the
server could follow only if it learned which points are reached at each
step.

\subsection{Homomorphic Comparison}
\label{sec:comparison}

The sign function is
\[
  \operatorname{sgn}(x) =
  \begin{cases}
    1 & \text{if } x > 0, \\
    0 & \text{if } x = 0, \\
    -1 & \text{if } x < 0.
  \end{cases}
\]
Every comparison in DBSCAN reduces to a sign. For two values $u$ and $v$,
we write
\[
  \operatorname{comp}(u,v) = \frac{1 + \operatorname{sgn}(v-u)}{2}
  = \begin{cases}
      1 & \text{if } u < v, \\
      \tfrac{1}{2} & \text{if } u = v, \\
      0 & \text{if } u > v,
    \end{cases}
\]
so that $\operatorname{comp}(u,v)$ is the outcome of the comparison
$u < v$ whenever $u \ne v$. The maximum follows from it as well:
\[
  \max(u,v) = u + (v-u) \cdot \operatorname{comp}(u,v).
\]
This expression also holds for $u = v$, since $v - u$ then vanishes.
The sign function is discontinuous at the origin, so no polynomial
represents it, and it can only be approximated on a domain that excludes
a neighborhood of the origin. We follow Lee et
al.~\cite{lee2021minimax}, who compose low-degree minimax polynomials
$p_1, \dots, p_k$ so that, for a precision parameter $\alpha$,
\begin{gather*}
  \bigl| (p_k \circ \cdots \circ p_1)(x) - \operatorname{sgn}(x) \bigr|
  \le 2^{1-\alpha} \\
  \text{for all } x \in [-1,1] \text{ with } |x| \ge 2^{-\alpha}.
\end{gather*}
Replacing $\operatorname{sgn}$ by this composition in
$\operatorname{comp}(u,v)$ then gives a value within $2^{-\alpha}$ of the
true outcome whenever $v - u \in [-1,1]$ and $|v - u| \ge 2^{-\alpha}$.
This is the comparison with precision $\alpha$
of~\cite{lee2021minimax}, and every comparison in
Section~\ref{sec:method} takes this form. A composition reaches this
accuracy at a far smaller multiplicative depth than a single polynomial of
the same accuracy. Even so, one evaluation spends a large part of the
level budget, so the ciphertext is bootstrapped between the components,
with the sign bootstrapping of Appendix~\ref{app:signboot}. For each
$\alpha$, we obtain the degrees of the components as
in~\cite{lee2021minimax}, taking into account that every component is
followed by this bootstrapping.

The approximation is applied to every slot at once, so one evaluation
decides $N/2$ comparisons. How a comparison outcome and a maximum are
built from it is shown in Section~\ref{sec:method}, where each is first
needed.

Two properties of this approximation are used throughout
Section~\ref{sec:method}. Its input must lie in $[-1,1]$, so every
quantity to be compared is scaled into that interval first. And the
interval $(-2^{-\alpha}, 2^{-\alpha})$, which we call the \emph{dead
zone}, carries no guarantee at all: an input that falls inside it may
evaluate to any value in $[-1,1]$. An algorithm built on this comparison
must therefore keep the quantities that decide its output away from the
dead zone, at a distance that $\alpha$ controls.
